\section{Introdução}
\label{sec:introducao}

\paragraph{}A simulação numérica consolidou-se como instrumento central da engenharia contemporânea, complementando abordagens teóricas e experimentais na investigação de fenômenos físicos complexos. No contexto da Engenharia Mecânica, a Dinâmica dos Fluidos Computacional (\textit{Computational Fluid Dynamics} -- CFD) e a modelagem de transferência de calor viabilizam a análise de problemas governados por Equações Diferenciais Parciais (EDPs), muitas vezes inviáveis por solução analítica direta ou por experimentação em escala de laboratório \cite{Maliska2004,Versteeg2007}. Essa capacidade preditiva tem papel relevante no dimensionamento, na otimização e na verificação de desempenho de sistemas térmicos e fluidodinâmicos.

\paragraph{}Apesar dos avanços metodológicos, o acesso a ambientes de simulação permanece limitado por barreiras econômicas e operacionais. Soluções comerciais amplamente utilizadas exigem licenciamento oneroso, infraestrutura computacional específica e configuração de ambiente com elevado custo de entrada, sobretudo em cenários de ensino e pesquisa inicial. Como consequência, a formação em métodos numéricos frequentemente fica condicionada à disponibilidade de laboratório e suporte técnico especializado, o que reduz a reprodutibilidade e a acessibilidade das atividades acadêmicas \cite{Anjos2013}.

\paragraph{}Neste contexto, este trabalho propõe o \textit{MinervaMesh}, uma plataforma de simulação numérica executada integralmente no navegador, com arquitetura \textit{client-side}. A solução combina Python científico com tecnologias web modernas, notadamente \textit{WebAssembly} e \textit{PyScript}, permitindo a execução local de rotinas de MEF e MDF sem instalação de dependências pelo usuário \cite{PyScript2023,Harris2020}. A contribuição principal não se restringe à implementação computacional, mas inclui a organização de um ambiente didático de engenharia para formulação, solução e análise de problemas térmicos e fluidodinâmicos.

\section{Objetivos}

\paragraph{}O objetivo geral deste trabalho é desenvolver e validar o \textbf{MinervaMesh}, uma plataforma computacional baseada na web para solução de EDPs de interesse em Engenharia Mecânica, operando integralmente no lado do cliente (\textit{client-side}).

\paragraph{}Os objetivos específicos são:
\begin{enumerate}
    \item Implementar métodos numéricos clássicos, com ênfase no Método dos Elementos Finitos (MEF) e uso complementar do Método das Diferenças Finitas (MDF), empregando bibliotecas científicas Python no ambiente web \cite{Harris2020}.
    \item Desenvolver uma interface interativa para pré-processamento, solução e pós-processamento, incluindo parametrização de malha e visualização de campos de interesse.
    \item Validar a plataforma, prioritariamente, por meio de casos com solução analítica conhecida e, em seguida, demonstrar sua aplicação em casos adicionais de engenharia, com foco em mecânica dos fluidos e transferência de calor.
    \item Consolidar um fluxo de trabalho reprodutível para problemas de:
    \begin{itemize}
        \item Escoamentos potenciais;
        \item Problemas de advecção-difusão (com estabilização SUPG) \cite{Brooks1982};
        \item Equações de Navier-Stokes para escoamentos laminares incompressíveis.
    \end{itemize}
\end{enumerate}

\section{Justificativa}

\paragraph{}A relevância do trabalho reside na articulação entre rigor numérico e acessibilidade computacional. Ao viabilizar simulações diretamente no navegador, reduz-se o atrito inicial de aprendizado e amplia-se o potencial de uso em disciplinas de graduação, monitorias e estudos independentes. Adicionalmente, a explicitação da formulação matemática e de sua implementação computacional fortalece o caráter acadêmico da proposta, permitindo que a plataforma seja utilizada tanto como ferramenta de simulação quanto como objeto de estudo em métodos numéricos \cite{Fish2007,Reddy2005}.

\section{Organização do Trabalho}

\paragraph{}Este trabalho está estruturado da seguinte forma:
\begin{itemize}
    \item O \textbf{Capítulo 2} (Revisão Bibliográfica) apresenta o histórico e o estado da arte dos temas que sustentam o trabalho, incluindo problemas térmicos, MDF, MEF e trabalhos correlatos.
    \item O \textbf{Capítulo 3} (Fundamentação Teórica) apresenta o embasamento matemático necessário para a discretização de EDPs via MEF.
    \item O \textbf{Capítulo 4} (Desenvolvimento e Arquitetura) detalha a implementação do \textit{MinervaMesh}, com ênfase na arquitetura \textit{client-side} e PyScript.
    \item O \textbf{Capítulo 5} (Resultados e Validação) apresenta a verificação numérica em casos com solução analítica e casos de aplicação.
    \item O \textbf{Capítulo 6} (Conclusão) resume as contribuições e propõe trabalhos futuros.
\end{itemize}
